\setchapterpreamble[u]{\margintoc}

To do.

Mention prediction markets as a source of distributed information and as a benchmark and use case of search.



Tying back to the central theme of search, the conclusion outlines what the next generation of information access may look like. We argue that \emph{omnimodal} search will become central as the set of exploitable sources of information widens with model capabilites (text, images, audio, video, code, structured data, 3D models).
Recent projects such as BrowseComp illustrate the growing impact of AI agents on real-world information seeking and highlight that robust access to complex information benefits from \emph{adaptive compute}. Crucially, these systems still benefit from strong ``first-layer'' retrieval primitives such as embedding-based search for efficient candidate selection.

The conclusion further reflects on the fact that verifying search correctness is not uniformly easy. Instead, search exhibits a \emph{jagged frontier} of capabilities: some tasks (e.g., factual verification grounded in accessible digital sources) are comparatively straightforward to check and should steadily improve \sidenote{Although as demonstrated by Vidore V3, complex queries remain very lengthy to verify even for human experts.}, whereas underspecified requests, open-ended synthesis, and value-laden questions remain harder both to solve and to verify. In this lens, verifier-side constraints shape what kinds of search problems will be reliably improved ``verifier's law''. (Mention the fact a lot of weak signals are provided by the user in a tyoical search system and can be exploited to understand whether a system has satisified a user request (are they asking follow ups, etc).

We further note that not all information is equally accessible: some sources remain costly to reach, poorly indexed, private, ephemeral... Exploiting such information at scale will therefore remain challenging in the coming years. Looking forward, we anticipate a shift from explicit querying toward more \emph{implicit} information access, where systems proactively surface relevant knowledge continuously in response to context, goals, and ongoing activity.

Finally, we argue that long-context modeling and information retrieval are closer than they may appear, and that bridging these literatures can be fruitful: both fields study how to select, compress, and condition on relevant evidence, finding needles of relevamnt information in a haystack of noise. Building conceptual and methodological links between them may be possible.

Beyond the technical horizon, we close by reflecting on what truly frictionless information access would mean, and on the societal implications of delegating discovery, synthesis, and decision support to increasingly agentic systems.